%% kapa: exam=1 round=23 year=2012 period=2 form=A subject=accounting subject_ko=회계학 seq=8 number=8
%% source: Q-Net 감정평가사 기출문제 (cst00309) artlSeq=5052580 file=raw/1008/5052580/2012년 제23회 감정평가사 제1차시험 기출문제(2교시 A형).pdf
%% conversion: kordoc 4.24.51 -> tools/gen1.py (auto); review=proofread-sonnet; flags=none
\begin{question}{8}
유형자산의 감가상각에 관한 설명으로 옳지 않은 것은?
\begin{choices}
  \item 유형자산을 구성하는 일부의 원가가 당해 유형자산의 전체원가에 비교하여 유의적이라면, 해당 유형자산을 감가상각할 때 그 부분은 별도로 구분하여 감가상각한다.
  \item 유형자산의 전체원가에 비교하여 해당 원가가 유의적이지 않은 부분도 별도로 분리하여 감가상각할 수 있다.
  \item 각 기간의 감가상각액은 다른 자산의 장부금액에 포함되는 경우가 아니라면 당기손익으로 인식한다.
  \item 유형자산의 잔존가치와 내용연수는 적어도 매 회계연도말에 재검토한다.
  \item 감가상각방법은 해당 자산의 공정가치 감소형태에 따라 선택한다.
\end{choices}
\end{question}
\answer{5}
